\chapter{Tutorial 1: Getting Set Up}
\begin{center}
\textit{January 3, 2018}
\end{center}

Welcome to tutorial 1! (This surprisingly happened even before the first day of
classes.)

\section{Course Infrastructure}
CS 246 course content is distributed through \code{git}, so we will need some
\code{git} commands; they are explained as we go through the course.

A Linux environment is preferred (and recommended) for all work. All course
material is available on the Linux student server, and assignments are tested
there, so code should always be checked on the server before submitting.

\section{Logging In to the Student Server}
\begin{lstlisting}[style=bash]
ssh username@linux.student.cs.uwaterloo.ca
# then enter the password for this account
\end{lstlisting}
If you do not have any clue what your \code{linux.student.cs} password was
from CS 136, it can be reset through the CS Computing Facility account web page.

This method works directly on:
\begin{enumerate}
  \item Linux,
  \item macOS,
  \item Bash on Ubuntu on Windows (Windows Subsystem for Linux; Windows 10
        Creators Update or later). Setup instructions:
        \url{https://docs.microsoft.com/en-us/windows/wsl/about}.
\end{enumerate}
Otherwise, download the \textbf{PuTTY} client (a free \code{.exe}) or
\textbf{SuperPuTTY}.

\subsection*{Copying files}
Files can be copied to and from the server with \code{scp}:
\begin{lstlisting}[style=bash]
scp a1q1.cc username@linux.student.cs.uwaterloo.ca:~/cs246/a1/
scp username@linux.student.cs.uwaterloo.ca:~/cs246/a1/out.txt .
\end{lstlisting}

\section{Course Directory and A0}
The course account has a directory (\code{/u/cs246/}) that holds the current
term's material; older terms' material is kept in older directories. A0 has been
posted (Piazza has hints about where A0 lives each term).

\begin{important}[Advice]
Start the assignments early.
\end{important}

\section*{Summary}
Use the Linux student server (via \code{ssh}/PuTTY), get course material with
\code{git}, and complete A0 (which is required) early.
